\chapter{Breast}

\label{ch:ch53_breast}

\section{What is breast imaging?}
Breast imaging uses mammography, ultrasound, and MRI to detect and characterize breast abnormalities. It is central to breast cancer screening, diagnosis, staging, and treatment monitoring.

\section{When is breast imaging ordered?}
\begin{itemize}
    \item \textbf{Screening:} Routine mammography in asymptomatic women (programmes vary by country; typically ages 50--70, biennial).
    \item \textbf{Symptomatic assessment:} A palpable lump, nipple discharge, skin changes, or pain.
    \item \textbf{High-risk screening:} Annual MRI for women with $>$20\% lifetime risk (BRCA carriers, prior chest radiation).
    \item \textbf{Pre-operative planning:} Assessing extent of known cancer, multifocality, and contralateral disease.
    \item \textbf{Post-treatment surveillance.}
\end{itemize}

\section{Mammography: 2D vs 3D (tomosynthesis)}

\begin{table}[h]
\centering
\caption{Mammography techniques}
\label{tab:mammo_techniques}
\begin{tabular}{@{}p{0.16\textwidth}p{0.34\textwidth}p{0.42\textwidth}@{}}
\hline
\textbf{Feature} & \textbf{2D digital} & \textbf{3D tomosynthesis} \\
\hline
Best for & General screening & Dense breasts, prior cancer \\
Detection rate & Baseline & +27--40\% over 2D \\
Dose & Standard & $<$2$\times$ 2D dose \\
Image acquisition & Single projection & Multiple low-dose projections \\
Recall rate & Higher & Lower (fewer false positives) \\
\hline
\end{tabular}
\end{table}

\section{Breast density}

\begin{table}[h]
\centering
\caption{ACR breast density categories}
\label{tab:density}
\begin{tabular}{clll}
\hline
\textbf{Cat.} & \textbf{Description} & \textbf{Composition} & \textbf{Risk} \\
\hline
A & Almost entirely fatty & $<$25\% glandular & Baseline \\
B & Scattered fibroglandular density & 25--50\% & 1.0--1.6$\times$ \\
C & Heterogeneously dense & 50--75\% & 1.7--2.1$\times$ \\
D & Extremely dense & $>$75\% & 2.1--4.0$\times$ \\
\hline
\end{tabular}
\end{table}

\begin{pearlbox}[Dense breasts]
Dense breast tissue (categories C and D) is found in approximately 40--50\% of women. It reduces mammographic sensitivity from $>$90\% in fatty breasts to as low as 30--48\% in extremely dense breasts while simultaneously increasing cancer risk. Women with dense breasts may benefit from supplemental screening ultrasound or MRI.
\end{pearlbox}

\section{Breast ultrasound}
Ultrasound is the primary adjunct to mammography:
\begin{itemize}
    \item Differentiates cysts from solid masses
    \item Guides percutaneous biopsies
    \item First-line investigation in women $<$30 years
    \item Evaluates focal pain, palpable abnormalities, and axillary lymph nodes
    \item No radiation---safe in pregnancy
\end{itemize}

\section{Breast MRI}
Breast MRI is the most sensitive imaging modality for breast cancer ($>$95\% sensitivity). Indications include:
\begin{itemize}
    \item High-risk screening (BRCA1/2, prior chest radiation, strong family history)
    \item Pre-operative assessment of extent of disease
    \item Assessment of neoadjuvant chemotherapy response
    \item Silicone implant integrity
    \item Occult primary (axillary metastases with negative mammogram/ultrasound)
\end{itemize}

\section{BI-RADS classification}

\begin{table}[h]
\centering
\caption{BI-RADS assessment categories}
\label{tab:birads}
\begin{tabular}{clll}
\hline
\textbf{Score} & \textbf{Term} & \textbf{Malignancy risk} & \textbf{Management} \\
\hline
0 & Incomplete & --- & Additional imaging needed \\
1 & Negative & 0\% & Routine screening \\
2 & Benign & 0\% & Routine screening \\
3 & Probably benign & $<$2\% & Short-interval follow-up (6 months) \\
4A & Low suspicion & 2--10\% & Tissue sampling recommended \\
4B & Moderate suspicion & 10--50\% & Tissue sampling recommended \\
4C & High suspicion & 50--95\% & Tissue sampling recommended \\
5 & Highly suggestive & $>$95\% & Tissue sampling / excision \\
6 & Known malignancy & --- & Staging before definitive treatment \\
\hline
\end{tabular}
\end{table}

\begin{pearlbox}[BI-RADS 3: the safety net]
A BI-RADS 3 assessment should only be used when the finding is genuinely probably benign. The lesion must be followed with imaging at 6, 12, and 24 months. If stable for 2 years, it can be downgraded to BI-RADS 2. Overuse of BI-RADS 3 delays cancer diagnosis.
\end{pearlbox}

\section{Screening vs diagnostic mammography}

\begin{table}[h]
\centering
\caption{Screening vs diagnostic mammography}
\label{tab:screen_vs_diag}
\begin{tabular}{@{}p{0.16\textwidth}p{0.34\textwidth}p{0.42\textwidth}@{}}
\hline
\textbf{Feature} & \textbf{Screening} & \textbf{Diagnostic} \\
\hline
BI-RADS & Should only assign 1 or 2 & Can assign 0--5 \\
Interpretation & Double reading + consensus & Single expert reader \\
Patient & Asymptomatic & Symptomatic or abnormal screening \\
Purpose & Detect early cancer & Characterise a finding \\
Views & Standard 2-view & Additional views (spot compression, magnification) \\
\hline
\end{tabular}
\end{table}

\section{Interventional breast procedures}
\begin{itemize}
    \item \textbf{Percutaneous biopsy:} Ultrasound-guided core needle biopsy is the gold standard for sampling suspicious masses. Vacuum-assisted biopsy is used for microcalcifications.
    \item \textbf{Preoperative localisation:} Wire or radioactive seed localisation marks non-palpable cancers for surgical excision.
    \item \textbf{Sentinel lymph node biopsy:} Not imaging per se, but guided by lymphoscintigraphy or blue dye.
    \item \textbf{Abscess drainage:} Ultrasound-guided aspiration of breast abscesses.
\end{itemize}

\section{Contraindications to screening mammography}
\begin{itemize}
    \item Absolute: Known BRCA or other high-penetrance mutation with $<$25 years
    \item Relative: Very young age (under 30) without high risk---dense tissue reduces utility
    \item Pregnancy: Avoid radiation; use ultrasound instead
    \item Active breast infection: Deferred until resolution
    \item Recent biopsy: Wait 1--2 weeks to avoid mistaking post-procedural change for pathology
\end{itemize}

\section{Reporting checklist}
A structured breast report should include:
\begin{enumerate}
    \item Indication and clinical history
    \item Breast density (ACR category)
    \item Asymmetries, masses (shape, margin, density, associated calcifications)
    \item Calcification morphology and distribution
    \item Architectural distortion
    \item Skin or nipple changes
    \item Axillary lymph nodes
    \item Comparison with prior imaging
    \item BI-RADS category
    \item Clear management recommendation
\end{enumerate}

\clearpage
